
% ============================================================
%  CHAPTER 7 -- System Performance
% ============================================================
\chapter{System Performance}
\label{ch:system-performance}

This chapter evaluates the performance of the five subsystems
comprising the gantry crane system: EERI I, EERI II, MCTR I,
MCTR II, and REII. The performance of each subsystem is assessed
against its design requirements using appropriate measurable
indicators and verification methods. The interactions between
subsystems are also considered to determine whether the integrated
system satisfies the overall design objectives.

\section{Performance Evaluation Methodology}
\label{sec:performance-methodology}

The performance of each subsystem shall be evaluated by comparing
its specified design requirements with the results obtained during
testing. Each performance indicator shall have a defined target,
a measurement method, and an acceptance criterion.

The following information shall be recorded for each indicator:
\begin{itemize}
    \item \textbf{Design requirement:} The required subsystem
    functionality or performance.
    \item \textbf{Performance indicator:} The measurable quantity
    used to evaluate the requirement.
    \item \textbf{Target value:} The value or operating range
    specified during the design process.
    \item \textbf{Measured result:} The value obtained during testing.
    \item \textbf{Assessment:} Whether the measured result satisfies
    the specified acceptance criterion.
\end{itemize}

\section{EERI I Performance}
\label{sec:eeri-i-performance}

This section evaluates the performance of EERI I with respect to
its assigned function within the gantry crane system. The indicators
selected for evaluation shall correspond to the subsystem's
functional and engineering requirements.

\begin{table}[H]
    \centering
    \caption{Performance indicators for EERI I.}
    \label{tab:eeri-i-performance}
    \renewcommand{\arraystretch}{1.4}
    \begin{tabularx}{\textwidth}{|X|X|X|X|}
        \hline
        \textbf{Performance indicator} &
        \textbf{Target value} &
        \textbf{Measured result} &
        \textbf{Assessment} \\
        \hline
        Primary functional performance &
        [Specified target] &
        [Measured result] &
        [Pass/Fail] \\
        \hline
        Operational accuracy or output quality &
        [Specified target] &
        [Measured result] &
        [Pass/Fail] \\
        \hline
        Response time or operating speed &
        [Specified target] &
        [Measured result] &
        [Pass/Fail] \\
        \hline
        Power consumption or operating efficiency,
        where applicable &
        [Specified target] &
        [Measured result] &
        [Pass/Fail] \\
        \hline
        Reliability under specified operating conditions &
        [Specified target] &
        [Measured result] &
        [Pass/Fail] \\
        \hline
    \end{tabularx}
\end{table}

The results shall be used to determine whether EERI I satisfies
its allocated requirements. Any deviations from the target values
shall be documented together with their probable causes and
effects on the overall system.

\section{EERI II Performance}
\label{sec:eeri-ii-performance}

The performance of EERI II shall be assessed according to its
role in the gantry crane system. Particular attention shall be
given to its principal function, the consistency of its output,
and its ability to operate within the specified design limits.

\begin{table}[H]
    \centering
    \caption{Performance indicators for EERI II.}
    \label{tab:eeri-ii-performance}
    \renewcommand{\arraystretch}{1.4}
    \begin{tabularx}{\textwidth}{|X|X|X|X|}
        \hline
        \textbf{Performance indicator} &
        \textbf{Target value} &
        \textbf{Measured result} &
        \textbf{Assessment} \\
        \hline
        Primary functional performance &
        [Specified target] &
        [Measured result] &
        [Pass/Fail] \\
        \hline
        Output accuracy or consistency &
        [Specified target] &
        [Measured result] &
        [Pass/Fail] \\
        \hline
        Response time or operating speed &
        [Specified target] &
        [Measured result] &
        [Pass/Fail] \\
        \hline
        Operating range or capacity, where applicable &
        [Specified target] &
        [Measured result] &
        [Pass/Fail] \\
        \hline
        Reliability under specified operating conditions &
        [Specified target] &
        [Measured result] &
        [Pass/Fail] \\
        \hline
    \end{tabularx}
\end{table}

The measured results shall establish whether EERI II performs
within its required operating range and whether its output is
suitable for the subsystem or subsystems with which it interacts.

\section{MCTR I Performance}
\label{sec:mctr-i-performance}

This section evaluates MCTR I in terms of its principal function
and its contribution to the mechanical and operational performance
of the gantry crane. Where applicable, the assessment shall include
load capacity, positional accuracy, motion, structural behaviour,
or other indicators relevant to the subsystem's actual design.

\begin{table}[H]
    \centering
    \caption{Performance indicators for MCTR I.}
    \label{tab:mctr-i-performance}
    \renewcommand{\arraystretch}{1.4}
    \begin{tabularx}{\textwidth}{|X|X|X|X|}
        \hline
        \textbf{Performance indicator} &
        \textbf{Target value} &
        \textbf{Measured result} &
        \textbf{Assessment} \\
        \hline
        Primary functional performance &
        [Specified target] &
        [Measured result] &
        [Pass/Fail] \\
        \hline
        Load capacity or operating capacity,
        where applicable &
        [Specified target] &
        [Measured result] &
        [Pass/Fail] \\
        \hline
        Positional accuracy or dimensional tolerance,
        where applicable &
        [Specified target] &
        [Measured result] &
        [Pass/Fail] \\
        \hline
        Operating speed or cycle time,
        where applicable &
        [Specified target] &
        [Measured result] &
        [Pass/Fail] \\
        \hline
        Mechanical integrity or stability,
        where applicable &
        [Specified criterion] &
        [Observed result] &
        [Pass/Fail] \\
        \hline
    \end{tabularx}
\end{table}

The results shall indicate whether MCTR I meets its required
functional and mechanical performance. Any observed deflection,
misalignment, excessive friction, instability, or other deviation
shall be documented where relevant to the subsystem.

\section{MCTR II Performance}
\label{sec:mctr-ii-performance}

MCTR II shall be evaluated against the requirements associated
with its assigned function. The assessment shall consider its
operational performance and any relevant mechanical, electrical,
or control-related characteristics.

\begin{table}[H]
    \centering
    \caption{Performance indicators for MCTR II.}
    \label{tab:mctr-ii-performance}
    \renewcommand{\arraystretch}{1.4}
    \begin{tabularx}{\textwidth}{|X|X|X|X|}
        \hline
        \textbf{Performance indicator} &
        \textbf{Target value} &
        \textbf{Measured result} &
        \textbf{Assessment} \\
        \hline
        Primary functional performance &
        [Specified target] &
        [Measured result] &
        [Pass/Fail] \\
        \hline
        Operating capacity or load handling,
        where applicable &
        [Specified target] &
        [Measured result] &
        [Pass/Fail] \\
        \hline
        Positioning accuracy or repeatability,
        where applicable &
        [Specified target] &
        [Measured result] &
        [Pass/Fail] \\
        \hline
        Response time or operating cycle time &
        [Specified target] &
        [Measured result] &
        [Pass/Fail] \\
        \hline
        Reliability and safe operation under
        specified conditions &
        [Specified criterion] &
        [Observed result] &
        [Pass/Fail] \\
        \hline
    \end{tabularx}
\end{table}

The results shall be compared with those of MCTR I where the
subsystems have related functions or operate together. This
comparison shall identify any performance differences that affect
the coordinated operation of the gantry crane.

\section{REII Performance}
\label{sec:reii-performance}

Final performance testing of the ESP32-S3-based REII
subsystem has not yet been completed. Earlier ESP32
radio measurements provide preliminary communication
results but do not establish integrated crane performance.
\cite{reii_measurements2026}

\begin{table}[H]
    \centering
    \caption{Performance indicators for REII.}
    \label{tab:reii-performance}
    \renewcommand{\arraystretch}{1.4}
    \begin{tabularx}{\textwidth}{|X|X|X|X|}
        \hline
        \textbf{Performance indicator} &
        \textbf{Target value} &
        \textbf{Measured result} &
        \textbf{Assessment} \\
        \hline
        Primary functional performance &
        [Specified target] &
        [Measured result] &
        [Pass/Fail] \\
        \hline
        Output accuracy or operational effectiveness &
        [Specified target] &
        [Measured result] &
        [Pass/Fail] \\
        \hline
        Response time or operating speed,
        where applicable &
        [Specified target] &
        [Measured result] &
        [Pass/Fail] \\
        \hline
        Operating capacity or efficiency,
        where applicable &
        [Specified target] &
        [Measured result] &
        [Pass/Fail] \\
        \hline
        Reliability or required safe response &
        [Specified criterion] &
        [Observed result] &
        [Pass/Fail] \\
        \hline
    \end{tabularx}
\end{table}

The results shall establish whether REII satisfies its allocated
requirements and whether its operation introduces any limitations
or dependencies affecting the performance of the complete system.

\section{Integrated System Performance}
\label{sec:integrated-performance}

The integrated performance assessment considers the interactions
between EERI I, EERI II, MCTR I, MCTR II, and REII. Individual
subsystem performance does not necessarily guarantee correct
system operation; therefore, the interfaces and coordinated
operating sequences shall also be evaluated.

\begin{table}[H]
    \centering
    \caption{Integrated system performance indicators.}
    \label{tab:integrated-performance}
    \renewcommand{\arraystretch}{1.4}
    \begin{tabularx}{\textwidth}{|X|X|X|X|}
        \hline
        \textbf{Performance indicator} &
        \textbf{Target value} &
        \textbf{Measured result} &
        \textbf{Assessment} \\
        \hline
        Correct operation of subsystem interfaces &
        [Specified criterion] &
        [Observed result] &
        [Pass/Fail] \\
        \hline
        Completion of the required operating sequence &
        [Specified criterion] &
        [Observed result] &
        [Pass/Fail] \\
        \hline
        Overall operating time or cycle time &
        [Specified target] &
        [Measured result] &
        [Pass/Fail] \\
        \hline
        System performance under the specified load
        or operating conditions &
        [Specified target] &
        [Measured result] &
        [Pass/Fail] \\
        \hline
        Fault handling and required safe-state behaviour &
        [Specified criterion] &
        [Observed result] &
        [Pass/Fail] \\
        \hline
    \end{tabularx}
\end{table}

Integration testing shall verify that signals, loads, power, and
other relevant quantities are transferred correctly across each
required interface. The complete operating sequence shall be
evaluated to identify delays, incompatibilities, or performance
limitations arising from interactions between subsystems.

\section{Summary of System Performance}
\label{sec:system-performance-summary}

The performance results for all five subsystems and the integrated
system shall be consolidated to assess compliance with the design
requirements.

\begin{table}[H]
    \centering
    \caption{Summary of subsystem performance.}
    \label{tab:system-performance-summary}
    \renewcommand{\arraystretch}{1.4}
    \begin{tabularx}{\textwidth}{|X|X|X|X|}
        \hline
        \textbf{Subsystem} &
        \textbf{Requirements assessed} &
        \textbf{Requirements satisfied} &
        \textbf{Outstanding issues} \\
        \hline
        EERI I & [Number or reference] &
        [Number or percentage] & [Details] \\
        \hline
        EERI II & [Number or reference] &
        [Number or percentage] & [Details] \\
        \hline
        MCTR I & [Number or reference] &
        [Number or percentage] & [Details] \\
        \hline
        MCTR II & [Number or reference] &
        [Number or percentage] & [Details] \\
        \hline
        REII & [Number or reference] &
        [Number or percentage] & [Details] \\
        \hline
        Integrated system & [Number or reference] &
        [Number or percentage] & [Details] \\
        \hline
    \end{tabularx}
\end{table}

The final assessment shall identify which requirements have been
satisfied, which remain outstanding, and whether corrective action
or further testing is required. The system shall be considered to
meet its performance objectives only when the relevant acceptance
criteria have been demonstrated by the recorded results.
